\section{Superglobal: \$\_GET, \$\_POST, dan \$\_SERVER}

Variabel \textbf{superglobal} adalah variabel bawaan PHP yang tersedia di seluruh scope --- baik di luar maupun di dalam fungsi --- tanpa perlu menggunakan kata kunci \code{global}. Superglobal berisi informasi penting tentang permintaan HTTP, data form, session, cookie, dan lingkungan server. Tiga superglobal yang paling sering digunakan dalam pemrosesan form adalah \code{\$\_GET}, \code{\$\_POST}, dan \code{\$\_SERVER}. Memahami cara kerja dan cara mengamankan superglobal ini adalah kunci untuk membangun aplikasi web yang fungsional dan aman \cite{phpManual}.

\code{\$\_GET} adalah array asosiatif yang berisi parameter dari query string URL. Ketika pengguna mengakses URL seperti \code{halaman.php?nama=Budi\&umur=20}, maka \code{\$\_GET["nama"]} berisi \code{"Budi"} dan \code{\$\_GET["umur"]} berisi \code{"20"}. Data GET terlihat di URL, sehingga cocok untuk parameter pencarian atau navigasi yang tidak sensitif. \code{\$\_POST} berisi data yang dikirim melalui form HTML dengan \code{method="post"}; data tidak terlihat di URL dan cocok untuk data sensitif atau jumlah data yang besar \cite{w3schoolsPhp}.

\begin{table}[ht]
\centering
\begin{tabular}{|P{2.5cm}|P{4.5cm}|P{4.5cm}|}
\hline
\textbf{Aspek} & \textbf{\$\_GET} & \textbf{\$\_POST} \\
\hline
Data dikirim via & Query string URL & Body HTTP request \\
\hline
Visibilitas & Terlihat di URL & Tidak terlihat di URL \\
\hline
Batas ukuran & Terbatas ($\sim$2048 karakter) & Hampir tidak terbatas \\
\hline
Bisa di-bookmark & Ya & Tidak \\
\hline
Keamanan & Kurang aman (terlihat) & Lebih aman (tersembunyi) \\
\hline
Kegunaan umum & Pencarian, filter, navigasi & Form login, registrasi, upload \\
\hline
Idempoten & Ya (seharusnya) & Tidak (boleh mengubah data) \\
\hline
\end{tabular}
\caption{Perbandingan \$\_GET dan \$\_POST}
\end{table}

\code{\$\_SERVER} adalah superglobal yang berisi informasi tentang server dan permintaan HTTP yang sedang masuk. Beberapa kunci (\textit{key}) yang paling sering digunakan antara lain \code{\$\_SERVER['REQUEST\_METHOD']} (metode HTTP: GET/POST), \code{\$\_SERVER['PHP\_SELF']} (path skrip yang sedang dieksekusi), \code{\$\_SERVER['HTTP\_HOST']} (nama host), dan \code{\$\_SERVER['REMOTE\_ADDR']} (alamat IP pengguna). Informasi ini sangat berguna untuk logika routing, logging, dan keamanan \cite{phpManual}.

\begin{webcode}[language=PHP, caption={Menggunakan superglobal GET, POST, dan SERVER}]
<?php
// Mengecek metode request
$method = $_SERVER["REQUEST_METHOD"];
echo "Metode: $method\n";

// Membaca data GET (dari URL: ?cari=php&hal=2)
$cari = $_GET["cari"] ?? "";
$halaman = (int) ($_GET["hal"] ?? 1);
echo "Pencarian: " . htmlspecialchars($cari) . "\n";

// Membaca data POST (dari form)
if ($method === "POST") {
    $email = $_POST["email"] ?? "";
    $pesan = $_POST["pesan"] ?? "";
    // Selalu sanitasi sebelum output!
    echo "Email: " . htmlspecialchars($email) . "\n";
}

// Informasi server
echo "Host: " . $_SERVER["HTTP_HOST"] . "\n";
echo "IP Pengguna: " . $_SERVER["REMOTE_ADDR"] . "\n";
echo "Skrip: " . $_SERVER["PHP_SELF"] . "\n";
?>
\end{webcode}

\begin{catatan}
\textbf{Prinsip Keamanan:} \textbf{Jangan pernah} langsung menggunakan data dari \code{\$\_GET}, \code{\$\_POST}, atau superglobal lainnya tanpa validasi dan sanitasi. Gunakan \code{htmlspecialchars()} saat menampilkan ke HTML, \code{filter\_input()} atau \code{filter\_var()} untuk validasi, dan \textit{prepared statements} untuk query database. Input pengguna harus \textbf{selalu} dianggap tidak tepercaya \cite{owaspXss}.
\end{catatan}

Selain \code{\$\_GET}, \code{\$\_POST}, dan \code{\$\_SERVER}, PHP memiliki beberapa superglobal lain yang akan dipelajari lebih lanjut di bab-bab berikutnya: \code{\$\_SESSION} (data session pengguna), \code{\$\_COOKIE} (data cookie), \code{\$\_FILES} (file yang diupload), dan \code{\$\_REQUEST} (gabungan GET, POST, dan COOKIE). Masing-masing memiliki fungsi spesifik dalam pengelolaan state dan data di aplikasi web \cite{phpRightWay}.

